\begin{pastexam}[label={px:ust-115-ee-ds-28}]{臺聯大電機類 115 學年度 第 28 題}
（單選；2 分）

Answer \textbf{Questions 28-30} according to the following C code.
\begin{minted}[linenos, numbersep=1.5mm, xleftmargin=2em, fontsize=\small]{c}
#include <stdio.h>
#include <stdlib.h>
int counter=0;
void max_heapify(int A[],int n,int i) {
    int largest=i;
    int left=2*i+1;
    int right=2*i+2;
    counter++;
    if (left<n && A[left]>A[i])
        largest=left;
    if (right<n && A[right]>A[largest])
        largest=right;
    if (largest!=i) {
        int temp=A[i];
        A[i]=A[largest];
        A[largest]=temp;
        max_heapify(A,n,largest);
    }
}
void build_heap(int A[],int n) {
    for (int i = (n/2)-1;i>=0;i--)
        max_heapify(A,n,i);
}
int main(){
    int A[]={4,1,3,2,16,9,10,14,8,7};
    int n = sizeof(A)/sizeof(A[0]);
    build_heap(A,n);
    for (int i=0;i<n;++i)
        printf("%d,",A[i]);
    printf("\n");
    printf("counter:%d\n",counter);
    return 0;
}
\end{minted}

What does the \texttt{\textbf{for-loop}} between \textbf{Lines 28 and 29} print?
\begin{enumerate}[label=(\Alph*)]
\item 4,1,3,2,16,9,10,14,8,7,
\item 16,14,10,8,7,9,3,2,4,1,
\item 1,2,3,4,7,8,9,10,14,16,
\item 16,14,10,8,7,3,9,1,4,2,
\item 16,14,10,9,8,7,4,3,2,1,
\end{enumerate}
\end{pastexam}

\begin{pastsoln}
答案是 (B)。

\texttt{build\_heap} 是由下而上建最大堆積（資結書第~8.2~節，虛擬碼~8.3），陣列從 $0$ 開始，
$A[i]$ 的孩子是 $A[2i+1]$、$A[2i+2]$。 $n = 10$，迴圈從 $i = 10/2 - 1 = 4$ 做到 $0$，每次把以
$A[i]$ 為根的子樹調成堆積（往下沉）。 追蹤（粗體是這一步交換過的格子）：
\begin{center}
\small
\begin{tabular}{l*{10}{c}}
\toprule
 & 0 & 1 & 2 & 3 & 4 & 5 & 6 & 7 & 8 & 9 \\
\midrule
開始 & 4 & 1 & 3 & 2 & 16 & 9 & 10 & 14 & 8 & 7 \\
$i = 4$ & 4 & 1 & 3 & 2 & 16 & 9 & 10 & 14 & 8 & 7 \\
$i = 3$ & 4 & 1 & 3 & \textbf{14} & 16 & 9 & 10 & \textbf{2} & 8 & 7 \\
$i = 2$ & 4 & 1 & \textbf{10} & 14 & 16 & 9 & \textbf{3} & 2 & 8 & 7 \\
$i = 1$ & 4 & \textbf{16} & 10 & 14 & \textbf{7} & 9 & 3 & 2 & 8 & \textbf{1} \\
$i = 0$ & \textbf{16} & \textbf{14} & 10 & \textbf{8} & 7 & 9 & 3 & 2 & \textbf{4} & 1 \\
\bottomrule
\end{tabular}
\end{center}
\begin{itemize}
\item $i = 4$：$A[4] = 16$ 只有左孩子 $A[9] = 7$，不用換。
\item $i = 3$：$A[3] = 2$ 的孩子 $14, 8$，和較大的 $14$（$A[7]$）交換；$A[7]$ 是葉子，停。
\item $i = 2$：$A[2] = 3$ 的孩子 $9, 10$，和 $10$（$A[6]$）交換，停。
\item $i = 1$：$A[1] = 1$ 的孩子 $14, 16$，和 $16$（$A[4]$）交換；$1$ 到了 $A[4]$，孩子 $A[9] = 7$，
再交換，$1$ 到 $A[9]$，停。
\item $i = 0$：$A[0] = 4$ 的孩子 $16, 10$，和 $16$ 交換；$4$ 到 $A[1]$，孩子 $14, 7$，和 $14$ 交換；
$4$ 到 $A[3]$，孩子 $2, 8$，和 $8$ 交換；$4$ 到 $A[8]$，停。
\end{itemize}
最後印出 \texttt{16,14,10,8,7,9,3,2,4,1,}。

其他選項：(A) 是原陣列（沒有建堆）；(C) 是由小到大排好的結果，建堆不會排序；(D) 前五個對了，但
$A[5], A[6]$ 的 $9, 3$ 被寫反、後面也不對；(E) 是由大到小排好的結果，那是堆積排序做完的樣子，
不是建堆的結果。
\end{pastsoln}
